% =====================================================================
%  DATA AND BACKGROUND  --  reframed, self-contained section
%
%  Replaces the current Section 3 of main.tex. Reframed around the paper's
%  actual analysis: a New Mexico ZIP-level border-exposure panel, a New
%  Mexico dispensary-level mechanism check, and a New Mexico vs. Colorado
%  county-level entry-regime comparison.
%
%  Figures/tables are embedded below. The three PNG files must be uploaded
%  to the project root (they live in "New Mexico project/figures and tables/"
%  and are also copied to "Latest try/reframe/figures/"):
%     fig_nm_setting.png
%     fig1_southwest_map.png          (replaces the old 5-state map)
%     fig_nm_spatial_inequality.png
%  New bibliography entries to append to bibliography.bib:
%     data_background_new_bib_entries.bib
%
%  Labels provided so pre-existing \ref{...}s elsewhere resolve:
%     sec:data, sec:samplecoverage, sec:sources, sec:summarystats,
%     sec:panelstructure, sec:timevarying, map, fig:dispensary_map,
%     tab:sample_construction (alias kept), tab:sumstats_nm, tab:sumstats_state
% =====================================================================

\section{Data and Background}
\label{sec:data}

The empirical analysis studies two states that share an important feature: each
operates a legal recreational cannabis market while sharing a long land border
with states that do not. New Mexico borders Texas, Oklahoma, and Utah, none of
which has an operating recreational market. Colorado borders Utah, Wyoming,
Nebraska, Kansas, and Oklahoma, none of which does either. Populations on the
nonlegal side of these borders are large: roughly 800,000 residents of Texas
live within 75 miles of the New Mexico line near Las Cruces alone. Dispensaries
in the border counties of both states therefore face a pool of potential
customers that interior dispensaries cannot serve as easily. Figure
\ref{fig:nm_setting} maps the New Mexico setting.

The two states differ in one respect that is central to the paper: how freely
firms can enter. New Mexico places no limit on the number of retail licenses and
issues them on an administrative timetable. Colorado sets no statewide license
cap but delegates licensing to local governments, and by 2019 roughly half of
its counties had banned recreational cannabis businesses altogether. Holding the
cross-border demand roughly fixed while varying the entry regime is what lets us
ask whether a border demand advantage is absorbed through additional entry, as in
a free-entry market, or through higher revenue at incumbent stores, as under
constrained entry.

\begin{figure}[htbp]
\centering
\includegraphics[width=0.92\textwidth]{fig_nm_setting.png}
\caption{The New Mexico setting. New Mexico operates a free-entry recreational
cannabis market and shares long borders with Texas, Oklahoma, and Utah, none of
which has a legal recreational market; its other neighbors, Colorado and
Arizona, do. Open red circles mark the centroids of the 92 New Mexico ZIP codes
classified as border-exposed (within 75 miles of a nonlegal-state line with at
least 1,000 nonlegal-state residents in reach); open blue circles mark the 276
interior ZIP codes. Gold points are the 1,192 geocoded dispensaries operating at
some point during 2022--2026.}
\label{fig:nm_setting}
\end{figure}

\subsection{Recreational cannabis licensing in New Mexico and Colorado}
\label{sec:timevarying}

\paragraph{New Mexico.} New Mexico legalized recreational cannabis through the
Cannabis Regulation Act, signed in April 2021 and effective June 29, 2021
\citep{NMCannabisRegulationAct}. Licensed retail sales began on April 1, 2022.
The Act is unusually explicit about entry: Section 26-2C-6(J) bars the Cannabis
Control Division from limiting the number of licensees or the number of licensed
premises a cannabis establishment may operate. Retail licenses are issued
administratively on an application basis, with the Division required to approve
or deny a complete application within a fixed window; there is no lottery, no
merit-scoring stage, no residency requirement, and no numerical cap. Local
governments may adopt time, place, and manner rules and density limits within
statutory bounds, but they cannot prohibit retail cannabis outright. In
practice, entry has been rapid and geographically dispersed: by the end of our
sample New Mexico had roughly 700 operating retailers spread across most of its
33 counties.

\paragraph{Colorado.} Colorado legalized recreational cannabis through Amendment
64, approved by voters in November 2012, with licensed retail sales beginning
January 1, 2014 \citep{COAmendment64}. The state imposes no numerical cap on
retail licenses, but Amendment 64 gives local governments full authority over
whether and how retail cannabis operates within their borders. A local
government may prohibit retail marijuana establishments entirely by ordinance or
ballot measure, or permit them subject to local caps on the number of stores and
restrictions on their location and hours. A retailer needs both a state license
from the Marijuana Enforcement Division and a local license. As of April 2019,
48.4 percent of Colorado counties banned all marijuana businesses
\citep{PayanBrownSong2021}, so entry in Colorado is best described as locally
capped even though no statewide cap exists.

\paragraph{Montana.} Montana, which enters only the pooled three-state
specification in Section~\ref{sec:regimeresults}, legalized recreational
cannabis through House Bill 701 (2021), with retail sales beginning January 1,
2022 \citep{MTHB701,MTBallotOptIn}. Entry is restricted in two ways: for the
first 18 months only existing medical-cannabis providers could obtain adult-use
licenses, and counties that voted against the 2020 legalization initiative must
approve adult-use sales by ballot before dispensaries can operate there. We
therefore code Montana as locally capped.

\subsection{Levels of analysis and sample coverage}
\label{sec:samplecoverage}
\label{sec:panelstructure}

The analysis works at three levels. The main border-exposure design
(Section~\ref{sec:borderexposure}) is a \emph{ZIP-code by quarter} panel for
New Mexico, covering all 368 populated ZIP codes over 2022Q2 through 2026Q3
(6,624 ZIP-quarters). The supply-side mechanism check
(Section~\ref{sec:mechanism}) is a \emph{dispensary by quarter} panel for New
Mexico, covering 1,192 geocoded dispensaries over the same window (10,739
dispensary-quarters). The cross-state entry-regime comparison
(Section~\ref{sec:crossstate}) is a \emph{county by quarter} panel for New
Mexico and Colorado, 2015Q1 (Colorado) or 2022Q2 (New Mexico) through 2025Q2.

The ZIP and dispensary panels for New Mexico are balanced: monthly
licensee-level reporting lets us observe every ZIP code and every dispensary in
every quarter, with a genuine zero for ZIP codes and quarters that record no
sales. The county panel is unbalanced. A county enters the estimation sample in
the first year it reports dispensary activity and can drop out later if its
revenue is suppressed for confidentiality; many of Colorado's rural
eastern-plains and Utah-border counties never enter at all. Colorado's revenue
source reports by fiscal year for some series, which we align to calendar year
for pooling. Table~\ref{tab:sumstats_state} reports the resulting county counts:
32 of New Mexico's 33 counties and 42 of Colorado's 64 counties appear in the
estimation sample.

\input{tab_sumstats_state.tex}

\subsection{Defining border and interior markets}

We classify markets as \emph{border-exposed} or \emph{interior} using geography
alone, so that treatment status is fixed before any dispensary chooses where to
locate. A ZIP code, dispensary, or county is border-exposed if two conditions
both hold: (i) its centroid (or, for a dispensary, its own address) lies within
75 miles of a state line shared with a state that has no operating recreational
cannabis market, and (ii) at least 1,000 residents of that nonlegal state live
within the same 75-mile reach. The population screen removes sparsely populated
stretches of state line, such as the empty desert along parts of the
New Mexico--Texas and Colorado--Utah borders, where proximity to the line
implies no meaningful cross-border demand. Neighbors with their own operating
recreational market count as legal: Arizona and Colorado for New Mexico
(Colorado from January 2022), and Arizona and, from April 2022, New Mexico for
Colorado. Markets near those borders are treated as interior.

Under this definition, 92 of New Mexico's 368 populated ZIP codes are
border-exposed and 276 are interior; 436 of its 1,192 dispensaries are
border-exposed. At the county level, the cross-state comparison in
Section~\ref{sec:crossstate} classifies 10 of New Mexico's 33 counties and 23 of
Colorado's 64 counties as border-exposed. Figure~\ref{fig:dispensary_map} maps
dispensary density and marks the counties that share a land border with a
nonlegal state (9 in New Mexico, 21 in Colorado); the two counts differ slightly
because the regression uses the 75-mile centroid rule while the map uses simple
geographic adjacency. Figure~\ref{fig:spatial_inequality} shows the
corresponding pattern within New Mexico at the ZIP level. Border and interior
ZIP codes have almost identical average populations (about 5,700 versus 5,600
residents), which is reassuring for a design that compares markets of similar
size differing mainly in their exposure to nonlegal-state demand.

\begin{figure}[htbp]
\centering
\includegraphics[width=\textwidth]{fig1_southwest_map.png}
\caption{Dispensary density and the nonlegal-state border, New Mexico and
Colorado, 2025. Shading is dispensaries per 1,000 residents in the 2025
estimation-sample cross-section; counties in white have no dispensary in that
cross-section. Red outlines mark the counties that share a land border with a
state that has no operating recreational cannabis market: 9 in New Mexico
(bordering Texas, Oklahoma, or Utah) and 21 in Colorado (bordering Utah,
Wyoming, Nebraska, Kansas, or Oklahoma). Counties whose only out-of-state
neighbors are legal-market states -- Arizona (green), or the shared
Colorado--New Mexico line, which are legal to each other after April 2022 -- are
interior. The county-level regressions in Section~\ref{sec:crossstate} instead
use a 75-mile centroid rule, which classifies 10 New Mexico and 23 Colorado
counties as border-exposed.}
\label{fig:dispensary_map}
\label{map}
\end{figure}

\begin{figure}[htbp]
\centering
\includegraphics[width=\textwidth]{fig_nm_spatial_inequality.png}
\caption{Spatial inequality in dispensary access, New Mexico. Panel (a): each
bubble is a ZIP-code centroid, sized by average dispensaries per 1,000
residents over 2022Q2--2026Q3 and colored by border status. Panel (b):
population-weighted dispensaries per 1,000 residents (solid bars, left axis) and
revenue per resident (hatched bars, right axis) for border and interior ZIP
codes, pooling all ZIP-quarters.}
\label{fig:spatial_inequality}
\end{figure}

\subsection{Data sources}
\label{sec:sources}

\paragraph{Cannabis sales.} New Mexico sales come from the Cannabis Control
Division's Cannabis Reporting Online Portal \citep{NMCCDReporting}, which
publishes monthly gross sales by individual licensee and by municipality. We
aggregate licensee-level sales to the ZIP code, the dispensary, and the county.
Colorado sales come from the Colorado Department of Revenue's monthly
county-level series \citep{COMarijuanaSalesByCounty}, spanning 2015 to the
present; we sum medical and recreational sales and keep only county-years with
all twelve months reported.

\paragraph{Dispensary locations and counts.} New Mexico dispensary addresses
come from the Cannabis Control Division's public list of licensed cannabis
establishments \citep{NMCCDLicensees}; we geocode each address with the U.S.
Census Bureau's batch geocoder and assign it to a ZIP code and county. Colorado
dispensary counts come from the Marijuana Enforcement Division's licensee lookup
\citep{COMedLicensees}, geocoded and assigned to counties by point-in-polygon.
For New Mexico, licensee-level reporting gives a true count of operating
dispensaries in every quarter; for Colorado, historical counts are reconstructed
from the geocoded license file.

\paragraph{Geography and population.} ZIP-code centroid coordinates and
population come from the SimpleMaps U.S. ZIP database
\citep{SimpleMapsZIP}, which is built on American Community Survey data. County
population is the U.S. Census Bureau's Population Estimates Program county total,
2023 vintage \citep{CensusPopEst2023}, held fixed across a county's panel years.
The population living within a given radius of a dispensary is computed from
2020 Census block populations and boundaries \citep{CensusBlocks2020}. State and
county boundaries are the Census Bureau's 2023 cartographic files
\citep{CensusCountyShapefile}.

\paragraph{Placebo sectors.} For the falsification tests in
Section~\ref{sec:falsification}, non-cannabis retail activity comes from the
New Mexico Taxation and Revenue Department's gross-receipts reports by location
and industry \citep{NMTRDGrossReceipts} and from the Census Bureau's ZIP Codes
Business Patterns \citep{CensusZBP}.

\subsection{Outcome variables and summary statistics}
\label{sec:summarystats}

We study three outcomes, each in natural logs. \textbf{Revenue per capita} is
total quarterly cannabis sales in a market divided by its population.
\textbf{Dispensary density} is the number of operating dispensaries per 1,000
residents. \textbf{Revenue per dispensary} is total quarterly sales divided by
the dispensary count, and is defined only for markets with at least one
operating dispensary. Revenue per capita and dispensary density capture the size
of the local market; revenue per dispensary captures how much any individual
store earns. Coefficients on log outcomes are reported throughout as percentage
effects, $(e^{\hat\beta}-1)\times 100$.

Tables~\ref{tab:sumstats_nm} and~\ref{tab:sumstats_state} report summary
statistics. Three patterns in the raw data preview the results. First, New
Mexico's border ZIP codes are much larger cannabis markets than its interior
ZIP codes: they average about six times the quarterly revenue per resident and
more than three times the dispensary density, even though the two groups have
nearly identical resident populations. Population-weighting narrows but does not
close the gap (Figure~\ref{fig:spatial_inequality}b): border ZIP codes still
support about 1.6 times the dispensary density and 1.7 times the revenue per
resident of interior ZIP codes. Second, that difference does not appear at the
individual-store level: border and interior dispensaries earn almost the same
average quarterly revenue (about \$222{,}000 versus \$202{,}000). Third,
Colorado's dispensaries earn roughly two and a half times the quarterly revenue
of New Mexico's, consistent with a more constrained entry environment spreading
demand across fewer stores.

\input{tab_sumstats_nm.tex}

\paragraph{Limitations.} The panels do not extend before New Mexico's April 2022
market opening, so the New Mexico analysis compares border and interior markets
after legalization rather than estimating the effect of legalization itself.
Population is measured at a single vintage and is not adjusted for the share of
residents old enough to purchase cannabis. Colorado's dispensary counts are
reconstructed from a current license file rather than a historical panel, which
introduces measurement error into Colorado's density and per-store outcomes but
not its revenue-per-capita outcome. Cross-border demand is proxied by nonlegal
population within a fixed radius; we do not observe individual consumers'
residence or travel.
